\documentclass[10pt,a4paper]{amsart}
\usepackage[margin=19mm]{geometry}
\usepackage{amsmath,amssymb,mathtools}
\usepackage[scr=rsfs,cal=euler]{mathalfa}
\usepackage{xfrac}
\usepackage[unicode,hidelinks,bookmarks=false]{hyperref}
\newcommand{\ie}{i.e.\ }
\newcommand{\cf}{cf.\ }
\newcommand{\resp}{respectively\ }
\newcommand{\Subsection}{\subsection}
\providecommand{\nobreakdash}{\nobreak\hbox{-}}
\setlength{\emergencystretch}{2em}
\multlinegap0pt
\usepackage{amssymb}
\usepackage{bm}
\usepackage{dsfont}
\usepackage{enumerate}
\usepackage{multirow}
\usepackage[normalem]{ulem}
\usepackage{mathtools}
\usepackage{booktabs}
\usepackage{tikz}
\newtheorem{theorem}{Theorem}
\newtheorem{lemma}{Lemma}
\newcommand{\Cb}{\mathbb{C}} %complexes
\newcommand{\Lcorner}{{\mathchoice{\mathrel{\makebox[7pt][c]{\rule{.4pt}{7.5pt}\rule{5pt}{.4pt}}}}{\mathrel{\makebox[7pt][c]{\rule{.4pt}{7.5pt}\rule{5pt}{.4pt}}}}{\mathrel{\makebox[5.5pt][c]{\rule{.4pt}{5.25pt}\rule{3.5pt}{.4pt}}}}{\mathrel{\makebox[4pt][c]{\rule{.4pt}{3.75pt}\rule{2.5pt}{.4pt}}}}}}
\newcommand{\Po}{\mathsf{P}}%sharp observable
\newcommand{\Qo}{\mathsf{Q}}%sharp observable
\newcommand{\Zb}{\mathbb{Z}} %integers
\newcommand{\id}{\mathds{1}} %identity operator
\newcommand{\ii}{\mathcal{J}}
\newcommand{\ip}[2]{\langle\,#1\,\vert\,#2\,\rangle} %inner product
\newcommand{\kb}[2]{\vert\,#1\,\rangle\langle\,#2\,\vert} %ketbra
\newcommand{\lc}[1]{\mathcal{L}(\Cb^{#1})} %linear operators on C^d
\newcommand{\no}[1]{\left\|\,#1\,\right\|} %norm
\newcommand{\phii}{\varphi} %basis
\newcommand{\ptr}[2]{{\rm tr}_{#1} #2} %partial trace
\newcommand{\rhoo}{\varrho} %states
\newcommand{\tr}[1]{{\rm tr}\left[#1\right]} %trace
\title{Postponing the Choice: Advantage of deferred measurements in quantum information processing}
\author{C. Carmeli, T. Heinosaari, A. Toigo}
\date{Source: arXiv:2502.19871v1 (CC BY 4.0)}
\begin{document}
\maketitle

\noindent\emph{Source and licence.} Postponing the Choice: Advantage of deferred measurements in quantum information processing. Version arXiv:2502.19871v1 (CC BY 4.0), distributed by arXiv under the Creative Commons Attribution 4.0 International licence, http://creativecommons.org/licenses/by/4.0/, as recorded in the arXiv licence metadata for this exact version and shown on the abstract page https://arxiv.org/abs/2502.19871v1. Full licence notice: \texttt{licenses/arxiv-2502.19871-CC-BY-4.0.txt}.\par\smallskip

\noindent\textit{Editorial scope.} Counted unit: the article's theorem thm:main (source0.tex lines 834--839). The article's complete original proof of it is delivered with it (source0.tex lines 1437--1529, one \texttt{proof} environment of 93 lines, which the article places away from the statement). No mathematics is written, repaired or re-derived: the statement and the proof are the article's own lines, in the article's own notation, and no retained line is substituted or re-flowed.  Retained uncounted as the source-local context the proof needs: lines 1040--1055; lines 1042--1047; lines 1344--1350; lines 1387--1389; lines 1396--1398; lines 1419--1421.  Nothing was omitted from any retained span, so the package carries no omission record.  No retained line contains a citation, so the wrapper carries no bibliography and no bibliography entry.  No retained line contains a figure environment, an image-inclusion command or drawing code, so no image is delivered and the package declares no asset dependency.  Editorial formatting only, in addition: an AMS \texttt{amsart} preamble that restates the article's own declarations and macros used by the retained lines, namely the environments \texttt{theorem}, \texttt{lemma} on separate counters, together with the standard packages those lines need.  The retained environments print in this document as Theorem 1, Lemma 1 in order of appearance, whereas the article prints the same items with the numbering of its own full document.  The complete native source of the article as distributed in the arXiv e-print for this exact version is bundled unchanged as \texttt{sources/173939/source0.tex}.
\begin{theorem}\label{thm:main_neg}
Suppose $(s,t)\in [m_1,1]\times [m_2,1]$. For any sharp meter $\Qo$, the meter $\Qo_s$ and the channel $I_t$ are compatible if and only if
\begin{equation}\label{eq:QI_comp_neg}
\begin{aligned}
& - 2\left[1 - t_d + \sqrt{(1 - d^2) \smash[b]{t_d}^2 + (d^2 - 2) \smash[b]{t_d} + 1}\right] \leq d^2 s \\
& \qquad\qquad \leq (d^2 - 2) (1 - t) + 2\sqrt{(1 - d^2) t^2 + (d^2 - 2) t + 1} \,,
\end{aligned}
\end{equation}
where
\begin{equation} %NOT CITED
t_d = \begin{cases}
\displaystyle \frac{d - 2}{2(d - 1)} & \displaystyle \text{ if } d \geq 3 \text{ and } t\leq\frac{d - 2}{2(d - 1)}\,, \\[0.3cm]
t & \text{ otherwise}\,.
\end{cases}
\end{equation}
\end{theorem}

\begin{equation}
\begin{aligned}
& - 2\left[1 - t_d + \sqrt{(1 - d^2) \smash[b]{t_d}^2 + (d^2 - 2) \smash[b]{t_d} + 1}\right] \leq d^2 s \\
& \qquad\qquad \leq (d^2 - 2) (1 - t) + 2\sqrt{(1 - d^2) t^2 + (d^2 - 2) t + 1} \,,
\end{aligned}
\end{equation}

\begin{lemma}[Symmetrization trick]\label{lem:trick}
Suppose $\Qo_s$ and $I_t$ are compatible. Then, there exists a joint instrument $\ii$ for $\Qo_s$ and $I_t$ such that
\begin{equation}\label{eq:Icov_def}
W(x,y)\ii(z,\rhoo)W(x,y)^* = \ii(x+z,W(x,y)\rhoo W(x,y)^*)
\end{equation}
for all $x,y,z$ and $\rhoo$.
\end{lemma}

\begin{equation}\label{eq:q}
q(z) = \tr{\vphantom{\bigg[}\Qo(z)\smash{\sum_x} S(x)}
\end{equation}

\begin{equation}\label{eq:ort_rel}
\sum_u \omega^{uv} = d\delta_{v,0}\,,
\end{equation}

\begin{equation}\label{eq:p}
p(x,z) = \tr{\Po(z)S(x)}
\end{equation}

\begin{theorem}\label{thm:main}
Suppose $(s,t)\in [0,1]\times [0,1]$. For any sharp meter $\Qo$, the noisy version $\Qo_s$ and the depolarizing channel $I_t$ are compatible if and only if
\begin{equation}\label{eq:QI_comp}
s+t-1\leq \frac{2}{d}\sqrt{(1-s)(1-t)}\, .
\end{equation}
\end{theorem}

\begin{proof}[Proof of Theorems \ref{thm:main} and \ref{thm:main_neg}]
The meter $\Qo_s$ is the convolution of the sharp meter $\Qo$ with the probability distribution $q_s$ given by
$$
q_s(z) = s\delta_{z,0} + \frac{1}{d}(1-s)
$$
for all $z$. In a similar way, the channel $I_t$ is the Weyl-covariant channel associated with the probability distribution $p_t$ defined as
$$
p_t(x,z) = t\delta_{x,0}\delta_{z,0} + \frac{1}{d^2}(1-t)\,.
$$
This follows from the fact that all states $\rhoo$ satisfy the equality
$$
\sum_{x,z} W(x,z)^*\rhoo W(x,z) = d\id\,,
$$
as it is easily checked by taking the matrix elements of both sides with respect to the standard basis and using \eqref{eq:ort_rel}. Comparing $q_s$ and $p_t$ with the probability distributions $q$ and $p$ of \eqref{eq:q} and \eqref{eq:p}, we see that the meter $\Qo_s$ and the channel $I_t$ admit a joint $W$-covariant instrument if and only if there exists a $\lc{d}$-valued vector measure $S$ on $\Zb_d$ which satisfies the two equations
\begin{equation}\tag{$\ast$}\label{eq:S_conditions_1}
\begin{aligned}
\tr{\vphantom{\bigg[}\Qo(z)\smash{\sum_x} S(x)} & = s\delta_{z,0} + \frac{1}{d}(1-s) \,,\\
\tr{\Po(z)S(x)} & = t\delta_{x,0}\delta_{z,0} + \frac{1}{d^2}(1-t)
\end{aligned}
\end{equation}
for all $x,z$. In turn, by Lemma \ref{lem:trick}, the existence of such $S$ is equivalent to the compatibility of $\Qo_s$ and $I_t$. In order to prove the theorems, for any fixed value of $t$, we will therefore determine the minimal and the maximal values $s_{\rm min}(t)$, $s_{\rm max}(t)$ of $s$ which allow the existence of a vector measure $S$ satisfying \eqref{eq:S_conditions_1}. Then, the  meter $\Qo_s$ and the channel $I_t$ are compatible for all $s\in [s_{\rm min}(t),s_{\rm max}(t)]$ by an easy convexity argument. To find $s_{\rm min}(t)$ and $s_{\rm max}(t)$, we fix $S$ as in \eqref{eq:S_conditions_1} and, for all $x$, we let $\xi(x)\in\Cb^d\otimes\Cb^d$ be a purification of $S(x)$, i.e., $S(x)=\ptr{\bar{1}}{\kb{\xi(x)}{\xi(x)}}$. Equations \eqref{eq:S_conditions_1} then rewrite
\begin{equation*}
\begin{aligned}
\sum_x \ip{\xi(x)}{(\Qo(z)\otimes\id)\xi(x)} & = s\delta_{z,0} + \frac{1}{d}(1-s) \,,\\
\ip{\xi(x)}{(\Po(z)\otimes\id)\xi(x)} & = t\delta_{x,0}\delta_{z,0} + \frac{1}{d^2}(1-t)\,.
\end{aligned}
\end{equation*}
By decomposing $\xi(x)$ as
$$
\xi(x) = \sum_z \psi_z\otimes\xi_z(x)\,,
$$
where $\xi_0(x),\ldots,\xi_{d-1}(x)$ are vectors of $\Cb^d$, the second equation becomes
$$
\no{\xi_z(x)}^2 = t\delta_{x,0}\delta_{z,0} + \frac{1}{d^2}(1-t) = \left(a_2(t)\delta_{x,0}\delta_{z,0} + \frac{1}{d} b(t)\right)^2
$$
for all $x,z$. Thus,
$$
\xi_z(x) = \left(a_2(t)\delta_{x,0}\delta_{z,0} + \frac{1}{d} b(t)\right) \eta_z(x)\,,
%\frac{1}{d}\sqrt{(d^2\delta_{x,0}\delta_{z,0} - 1) t + 1}\,\eta_z(x)
$$
where $\eta_z(x)$ is a unit vector of $\Cb^d$. By using the relation $\ip{\phii_u}{\psi_v} = \omega^{uv}/\sqrt{d}$, the first equation then yields
\begin{align*}
s\delta_{z,0} + \frac{1}{d}(1-s) & = \sum_{x,u,v} \ip{\phii_z}{\psi_u} \ip{\psi_v}{\phii_z} \ip{\xi_v(x)}{\xi_u(x)} \\
& = \frac{1}{d} \sum_x \bigg\|\sum_u \omega^{uz} \xi_u(x)\bigg\|^2 \\
& = \frac{1}{d} \sum_x \bigg\|\sum_u \omega^{uz} \left(a_2(t)\delta_{x,0}\delta_{u,0} + \frac{1}{d} b(t)\right) \eta_u(x)\bigg\|^2\,,
\end{align*}
which gives
$$
s = \frac{1}{d-1} \bigg[ \sum_x \bigg\|\sum_u \left(a_2(t)\delta_{x,0}\delta_{u,0} + \frac{1}{d} b(t)\right) \eta_u(x)\bigg\|^2 - 1 \bigg]
$$
in the particular case with $z=0$. By triangular inequality,
\begin{align*}
\bigg\|\sum_u \left(a_2(t)\delta_{x,0}\delta_{u,0} + \frac{1}{d} b(t)\right) \eta_u(x)\bigg\| & \leq a_2(t)\delta_{x,0} + b(t)
\end{align*}
and therefore
\begin{align*}
s & \leq \frac{1}{d-1} \big[ \big(a_2(t) + b(t)\big)^2 + (d-1)b(t)^2 - 1 \big] \\
& = \frac{1}{d^2}\left[(d^2 - 2) (1 - t) + 2\sqrt{(1 - d^2) t^2 + (d^2 - 2) t + 1}\right]\,.
\end{align*}
For $s$ attaining the above equality, a joint instrument for $\Qo_s$ and $I_t$ actually exists, namely, the instrument $\tilde{\ii}^{(t)}$. It follows that
$$
s_{\rm max}(t) = \frac{1}{d^2}\left[(d^2 - 2) (1 - t) + 2\sqrt{(1 - d^2) t^2 + (d^2 - 2) t + 1}\right]\,.
$$
This proves the second inequality of \eqref{eq:QI_comp_neg}. On the other hand, we have
\begin{align*}
& \sum_x \bigg\|\sum_u \left(a_2(t)\delta_{x,0}\delta_{u,0} + \frac{1}{d} b(t)\right) \eta_u(x)\bigg\|^2 \geq \\
& \qquad\qquad \geq \bigg\|\sum_u \left(a_2(t)\delta_{u,0} + \frac{1}{d} b(t)\right) \eta_u(0)\bigg\|^2 \geq \\
& \qquad\qquad \geq \bigg[\left(a_2(t) + \frac{1}{d} b(t)\right) - \frac{1}{d} b(t) \bigg\|\sum_u (1-\delta_{u,0}) \eta_u(0)\bigg\|\bigg]^2 \,,
\end{align*}
where the last relation follows from another application of the triangular inequality. By the same reason,
$$
\bigg\|\sum_u (1-\delta_{u,0}) \eta_u(0)\bigg\| \leq d-1 \,.
$$
If $t\geq (d-2)/[2(d-1)]$, we have
$$
a_2(t) + \frac{1}{d} b(t) \geq \frac{d-1}{d} b(t) \,,
$$
which implies
\begin{align*}
& \bigg[\left(a_2(t) + \frac{1}{d} b(t)\right) - \frac{1}{d} b(t) \bigg\|\sum_u (1-\delta_{u,0}) \eta_u(0)\bigg\|\bigg]^2  \\
& \qquad\qquad\qquad\qquad\qquad\qquad \geq \bigg[\left(a_2(t) + \frac{1}{d} b(t)\right) - \frac{d-1}{d} b(t) \bigg]^2 \,.
\end{align*}
In the case with $d=2$ the last inequality is true also for $t<(d-2)/[2(d-1)]$ and it is actually an equality. It follows that, for $d=2$ or $t\geq (d-2)/[2(d-1)]$,
\begin{align*}
s & \geq \frac{1}{d-1} \bigg\{ \bigg[\left(a_2(t) + \frac{1}{d} b(t)\right) - \frac{d-1}{d} b(t) \bigg]^2 - 1 \bigg] \\
& = - \frac{2}{d^2}\left[1 - t + \sqrt{(1 - d^2) t^2 + (d^2 - 2) t + 1}\right]\,.
\end{align*}
We can find a joint instrument for $\Qo_s$ and $I_t$ also when $s$ attains the last equality, namely, the instrument $\tilde{\ii}_-^{(t)}$. As a consequence, if we still assume $d=2$ or $t\geq (d-2)/[2(d-1)]$, then
$$
s_{\rm min}(t) = - \frac{2}{d^2}\left[1 - t + \sqrt{(1 - d^2) t^2 + (d^2 - 2) t + 1}\right] \,.
$$
In particular, $s_{\rm min}\big((d-2)/[2(d-1)]\big) = m_1$. When $d\geq 3$, we have also $s_{\rm min}(m_2) = m_1$ because $\Qo_{m_1}$ and $I_{m_2}$ admit the joint instrument $\ii_\Lcorner$, hence $s_{\rm min}(t) = m_1$ for all $t\in\big[m_2,(d-2)/[2(d-1)]\big]$ by the usual convexity argument. The first inequality of \eqref{eq:QI_comp_neg} then follows by combining the two cases with $d=2$ and $d\geq 3$. 
\end{proof}

\end{document}
